\section{Выводы}

Основное преимущество суффиксного массива перед наивными алгоритмами поиска
заключается в том, что после однократного построения массива поиск любого
образца выполняется за время $O(m \log n)$, где $m$ — длина образца,
$n$ — длина текста. Это делает его особенно полезным
в задачах, где текст не изменяется, а запросов на поиск много.

Проведённое тестирование производительности наглядно продемонстрировало сильные и
слабые стороны обоих подходов. Алгоритм КМП не требует предобработки текста и выполняет
поиск за время $O(n + m)$, что делает его эффективным для однократного поиска или в случаях,
когда текст часто меняется. Суффиксный массив, напротив, требует значительных затрат на построение,
но затем обеспечивает практически мгновенный поиск (доли миллисекунды для текста длиной в миллион символов).

Особенно показательным оказался тот факт, что наивное построение суффиксного массива
через стандартную сортировку имеет квадратичную сложность в худшем случае ($O(n^2 \log n)$), что
подтвердилось тестом на тексте из 1\,000\,000 символов — построение заняло более двух минут.
Это подчёркивает важность использования эффективных алгоритмов построения суффиксного массива, например,
построение через суффиксное дерево (сложность $O(n)$).

Таким образом, данная лабораторная работа позволила не только изучить теоретические основы
суффиксных массивов, но и на практике убедиться в компромиссе между временем предобработки и
временем поиска, что является фундаментальной концепцией в анализе алгоритмов и структур данных.

\pagebreak